\subsection{Worked example: bounded eventuality with an exclusion window}
\label{subsec:worked-translation}
Consider
\[
  f=(\Eventually_{\le5}p)\land(\Always_{<2}\neg p).
\]
It requires $p$ to be false at the origin and until age two, then true by age
five. Figure~\ref{fig:translation-reference-tba} shows a hand-simplified
equivalent one-clock TBA, not the syntax-directed output, which assigns
separate clocks to the two distinct timed subformulas. The steps below show
the clocked-LTL translation and timed reinterpretation.

\begin{figure}[t]
\centering
\resizebox{0.88\columnwidth}{!}{\begin{tikzpicture}[node distance=22mm]
  \node[fmstate] (q0) {$q_{\mathit{init}}$};
  \node[fmstate,right=of q0] (q1) {$q_{\mathit{wait}}$};
  \node[fmaccept,right=of q1] (q2) {$q_{\mathit{acc}}$};
  \draw[fmedge] ([xshift=-8mm]q0.west) -- (q0.west);
  \draw[fmedge] (q0) -- node[fmlabel,above] {$\neg p\;/\;\mathsf{rst}(x)$} (q1);
  \draw[fmedge] (q1) edge[loop above] node[fmlabel] {$x\le5,\;\neg p$} (q1);
  \draw[fmedge] (q1) -- node[fmlabel,above] {$2\le x\le5,\;p$} (q2);
  \draw[fmedge] (q2) edge[loop above] node[fmlabel] {$\top$} (q2);
\end{tikzpicture}
}
\caption{Hand-simplified equivalent event-guarded TBA for $f$. The initial
transition requires $\neg p$ and resets $x$; in the waiting location,
$x\le5$, and acceptance also requires $x\ge2$. Thus $p$ first holds at age
in $[2,5]$.}
\label{fig:translation-reference-tba}
\end{figure}

\begin{enumerate}[label=Step~\arabic*:,leftmargin=*]
\item The formula is in negation normal form. Expanding the derived operators
and applying the timed-operator identities gives
\[
\begin{aligned}
f
 &= (\top\Until_{\le5}p)\land(\bot\Release_{<2}\neg p)\\
 &= \bigl(p\lor(\top\land\top\hUntil_{\le5}p)\bigr)\\
 &\qquad\land\bigl(\neg p\land(\bot\lor\bot\hRelease_{<2}\neg p)\bigr).
\end{aligned}
\]
These current-position cases follow Eqs.~\eqref{eq:derive-ule} and
\eqref{eq:derive-rle}.

\item The two distinct primitive timed subformulas receive clocks $x$ and
$y$, respectively. Their clocked-LTL translation is
\begin{align*}
\Phi_U={}&p\lor\bigl(\top\land\Next\Bigl(
  ((x\le5)\land\Unch(x)\land\top)\Until\\[-1mm]
&\qquad((x\le5)\land p)\Bigr)\bigr),\\[-1mm]
\Phi_R={}&\neg p\land\bigl(\bot\lor(\Reset(y)\land\\[-1mm]
&\qquad\Next(\neg p\WeakUntil((y\ge2)\lor(\bot\land\neg p))))\bigr),\\[-1mm]
\T(f)={}&\Phi_U\land\Phi_R.
\end{align*}
The Until preserves $x$; the Release test $y\ge2$ complements the source
bound $<2$, allowing $p$ from age two.

\item A Spot backend output for $\T(f)$ is shown on the left of
Fig.~\ref{fig:translation-steps45}. Clock guards, reset and preservation
markers, and action atoms are independent propositions to the backend. In
particular, $p=\mathsf{LAct}(p)$ and $\neg p=\mathsf{LNAct}(p)$ are distinct
atoms, while $!p$ is the negated literal of $\mathsf{LAct}(p)$; a backend cube
may contain both $p$ and $\neg p$. Timed reinterpretation imposes action
consistency.

\item Relaxation removes negative literals over extended clock atoms. Reset
completion resets $x$ unless $\Unch(x)$ is present and resets $y$ exactly when
$\Reset(y)$ is present. The right side of Fig.~\ref{fig:translation-steps45}
shows the transitions after marker removal; later optimization removes
unsatisfiable action cubes and simplifies the automaton.
\end{enumerate}

\begin{figure*}[!t]
\centering
\resizebox{\textwidth}{!}{\begin{tikzpicture}[
  x=1cm,y=1cm,
  lab/.style={font=\small,align=center,fill=white,inner sep=1pt}
]
% Left: Spot 2.16 output; state IDs match examples/spot/worked_translation.dot.
\begin{scope}[xshift=0cm]
  \node[fmnote,font=\normalsize\bfseries] at (2.4,6.0) {Spot output};
  \node[fmstate] (l0) at (0.0,4.8) {$2$};
  \node[fmstate] (l1) at (0.0,2.4) {$3$};
  \node[fmaccept] (l2) at (4.8,4.8) {$1$};
  \node[fmstate] (l3) at (4.8,2.4) {$4$};
  \node[fmaccept] (l4) at (2.4,0.0) {$0$};
  \draw[fmedge] (-1.0,4.8) -- (l0);
  \draw[fmedge] (l0) -- (l1);
  \node[lab,anchor=east] at (-0.18,4.05) {$\begin{array}{r}!p\land\neg p\\{}\land\,\mathsf{rst}(y)\end{array}$};
  \draw[fmedge] (l0) -- node[lab,above,pos=.32,yshift=2pt] {$p\land\neg p\land\mathsf{rst}(y)$} (l2);
  \draw[fmedge] (l1) edge[loop left,looseness=5.5] node[lab,left,xshift=-3pt] {$\begin{array}{r}\mathsf{unch}(x)\land x\le5\\{}\land\,!p\land\neg p\\{}\land\,!(y\ge2)\end{array}$} (l1);
  \draw[fmedge] (l1) -- (l2);
  \node[lab] at (2.4,4.35) {$\begin{array}{c}x\le5\land p\land\neg p\\{}\land\,!(y\ge2)\end{array}$};
  \draw[fmedge] (l1) -- (l3);
  \node[lab] at (2.4,1.88) {$\begin{array}{c}\mathsf{unch}(x)\land x\le5\\{}\land\,!p\land y\ge2\end{array}$};
  \draw[fmedge] (l1) -- (l4);
  \node[lab] at (0.75,0.88) {$\begin{array}{r}x\le5\land p\\{}\land y\ge2\end{array}$};
  \draw[fmedge] (l2) edge[loop right,looseness=5.5] node[lab,right,xshift=3pt] {$\begin{array}{l}\neg p\\{}\land\,!(y\ge2)\end{array}$} (l2);
  \draw[fmedge] (l2) to[out=-15,in=10,looseness=1.3] node[lab,right,pos=.42,xshift=4pt] {$y\ge2$} (l4);
  \draw[fmedge] (l3) edge[loop right,looseness=5.5] node[lab,right,xshift=8pt] {$\begin{array}{l}\mathsf{unch}(x)\\{}\land x\le5\land\,!p\end{array}$} (l3);
  \draw[fmedge] (l3) -- (l4);
  \node[lab] at (4.05,1.32) {$x\le5\land p$};
  \draw[fmedge] (l4) edge[loop below,looseness=5] node[lab,below] {$\top$} (l4);
\end{scope}

% Right: relaxation and reset completion, with the same Spot state IDs.
\begin{scope}[xshift=9.6cm]
  \node[fmnote,font=\normalsize\bfseries] at (2.4,6.0) {Reset-completed TBA};
  \node[fmstate] (r0) at (0.0,4.8) {$2$};
  \node[fmstate] (r1) at (0.0,2.4) {$3$};
  \node[fmaccept] (r2) at (4.8,4.8) {$1$};
  \node[fmstate] (r3) at (4.8,2.4) {$4$};
  \node[fmaccept] (r4) at (2.4,0.0) {$0$};
  \draw[fmedge] (-1.0,4.8) -- (r0);
  \draw[fmedge] (r0) -- (r1);
  \node[lab,anchor=east] at (-0.18,4.05) {$\begin{array}{r}\neg p\\/\;\mathsf{rst}(x),\mathsf{rst}(y)\end{array}$};
  \draw[fmedge] (r0) -- node[lab,above,pos=.32,yshift=2pt] {$p\land\neg p\;/\;\mathsf{rst}(x),\mathsf{rst}(y)$} (r2);
  \draw[fmedge] (r1) edge[loop left,looseness=5.5] node[lab,left,xshift=-3pt] {$x\le5,\;\neg p$} (r1);
  \draw[fmedge] (r1) -- (r2);
  \node[lab] at (2.4,4.35) {$x\le5,\;p\land\neg p\;/\;\mathsf{rst}(x)$};
  \draw[fmedge] (r1) -- (r3);
  \node[lab] at (2.4,1.88) {$x\le5\land y\ge2,\;\neg p$};
  \draw[fmedge] (r1) -- (r4);
  \node[lab] at (0.75,0.88) {$\begin{array}{r}x\le5\land y\ge2,\;p\\/\;\mathsf{rst}(x)\end{array}$};
  \draw[fmedge] (r2) edge[loop right,looseness=5.5] node[lab,right,xshift=3pt] {$\begin{array}{l}\neg p\\/\;\mathsf{rst}(x)\end{array}$} (r2);
  \draw[fmedge] (r2) to[out=-15,in=10,looseness=1.3] node[lab,right,pos=.5,xshift=8pt] {$\begin{array}{l}y\ge2,\;\top\\/\;\mathsf{rst}(x)\end{array}$} (r4);
  \draw[fmedge] (r3) edge[loop right,looseness=5.5] node[lab,right,xshift=8pt] {$x\le5,\;\neg p$} (r3);
  \draw[fmedge] (r3) -- (r4);
  \node[lab] at (4.05,1.32) {$\begin{array}{l}x\le5,\;p\\/\;\mathsf{rst}(x)\end{array}$};
  \draw[fmedge] (r4) edge[loop below,looseness=5] node[lab,below] {$\top\;/\;\mathsf{rst}(x)$} (r4);
\end{scope}
\end{tikzpicture}
}
\caption{Spot 2.16 output (redrawn, left) and hand-derived reset-completed
TBA (right; same state IDs). Raw input and DOT are in
\texttt{examples/spot/}. Right-hand edges show guards, actions, and resets;
$\mathsf{rst}(c)$ and $\mathsf{unch}(c)$ mark reset and preservation.}
\label{fig:translation-steps45}
\end{figure*}
